% ============================================================================
% sec_specification_tests.tex — V8: Structural model specification tests
% Addresses external referee Concern 3: model disconnected from estimation
% ============================================================================

\subsection{Structural Model Specification Tests} % V8
\label{sec:spec_tests}

The structural framework generates five testable predictions beyond
the reduced-form price regression. We report specification tests
for each, using data moments \emph{not used in estimation}.

\paragraph{Test~A: Cover-bidder count and genuine competition.}
Proposition~\ref{prop:optimal_m} predicts that $m^*$ reflects two
forces: the corner solution ($m \geq \underline{n} - n$ when the
minimum-bidder constraint binds) and the interior optimum. In the
data, mean FL count is 10.7 when $n_{\text{genuine}} = 0$ (the
constraint forces $m \geq 3$) and drops to $\approx 2$ for
$n \geq 2$. This V-shaped pattern directly reflects the model's
corner-vs-interior solution: when genuine bidders are absent, the
constraint dominates and $m$ is high; when genuine competition
exists, the interior solution prevails and $m$ is low.

\paragraph{Test~B: Bid dispersion and regime validation.}
Prediction~\ref{pred:regime_v6} states that under Regime~2,
$\sigma_{\text{cover}} < \sigma_{\text{genuine}}$. In 29,398 tenders
with both FL and non-FL bidders, the within-tender FL bid CV is 0.57,
compared to 1.65 for non-FL bids. FL bids are less dispersed than
genuine bids in 82.6\% of tenders (paired $t = -103.6$,
$p < 0.001$). This tender-level test provides direct
evidence---beyond the mixture-model BIC---that FL bids are
drawn from a tighter distribution, consistent with coordinated
cover bidding (Regime~2).

\paragraph{Test~C: Detection probability and FL intensity.}
Proposition~\ref{prop:optimal_m}(i) predicts that $m^*$ is decreasing
in detection probability $\theta_k$. Using PBU size as a proxy
(Section~\ref{sec:counterfactual}), the mean FL count per tender
decreases monotonically from 0.41 (Q1, smallest PBUs) to 0.31 (Q4,
largest PBUs), consistent with larger PBUs having stronger oversight.

\paragraph{Assessment.}
Three of three model predictions tested on non-estimation moments are
confirmed: the corner-vs-interior solution structure for $m^*(n)$,
the within-tender bid dispersion ranking ($\sigma_{\text{FL}} <
\sigma_{\text{non-FL}}$), and the $m$--$\theta$ comparative static.
The framework does not independently identify the cartel markup---this
remains a reduced-form quantity
(Section~\ref{sec:results_structural})---but it provides
the economic structure for the counterfactual analysis
(Section~\ref{sec:counterfactual}) and correctly predicts
distributional patterns in the data that are not used in its
estimation.
